%=============================================================================
% WAVE 4, ITERATION 20: P10 (Computing / mochi_class Implementation)
% Agent 9: NONLINEAR mu-FUNCTION CRISIS --- Can mochi_class Handle DFD?
%
% Model: Opus 4.6
% Date: 20 March 2026
%
% INHERITED STATE (from i19):
%   - DFD maps to Horndeski with alpha_T=alpha_M=alpha_B=0, alpha_K~0.014
%   - Minimal mochi_class plan: 150 lines C, 1 developer-week
%   - FUNDAMENTAL DISAGREEMENT:
%       Agent 4 (P4): G_eff = G_N on all sub-Hubble scales (linearized)
%       Agent 12 (Cosmo): G_eff >> G_N (deep-MOND, nonlinear mu)
%   - Cosmo agent showed delta_b ~ a^2 growth in deep-MOND, rescuing P(k)
%   - P4 agent showed strict QSA + linearization gives G_eff = G_N exactly
%   - Deep-MOND regime involves nonlinear mu, which standard Horndeski
%     parametrization may NOT capture
%
% MANDATE (i20):
%   1. Can existing mochi_class handle DFD? Analyze the linear/nonlinear gap
%   2. What modifications to CLASS are needed for nonlinear mu?
%   3. Write MINIMAL code modifications (C, targeting perturbations.c)
%   4. Estimate P(k), C_l^TT, sigma_8 under both scenarios
%   5. WebSearch: MOND-like perturbation theory in Boltzmann codes
%=============================================================================

\documentclass[11pt]{article}
\usepackage{amsmath,amssymb,amsthm}
\usepackage[margin=1in]{geometry}
\usepackage{booktabs}
\usepackage{xcolor}
\usepackage{tcolorbox}
\usepackage{listings}
\usepackage{hyperref}
\usepackage{enumitem}

\newtheorem{theorem}{Theorem}[section]
\newtheorem{proposition}[theorem]{Proposition}
\newtheorem{lemma}[theorem]{Lemma}
\newtheorem{finding}[theorem]{Finding}
\theoremstyle{definition}
\newtheorem{definition}[theorem]{Definition}
\theoremstyle{remark}
\newtheorem{remark}[theorem]{Remark}

\newcommand{\ee}[1]{\times 10^{#1}}
\newcommand{\Geff}{G_{\rm eff}}
\newcommand{\kmu}{k_\mu}
\newcommand{\astar}{a_\star}
\newcommand{\LCDM}{$\Lambda$CDM}
\newcommand{\DFD}{\textsc{dfd}}
\newcommand{\dpsi}{\delta\psi}
\newcommand{\aM}{a_0}

\lstset{
  language=C,
  basicstyle=\ttfamily\small,
  keywordstyle=\color{blue},
  commentstyle=\color{green!50!black},
  stringstyle=\color{red!70!black},
  numbers=left,
  numberstyle=\tiny\color{gray},
  frame=single,
  breaklines=true,
  captionpos=b,
  showstringspaces=false
}

\tcbuselibrary{skins,breakable}

\title{\textbf{Wave 4 Iteration 20: Cross-Reference Update}\\[6pt]
Patent 10 --- The Nonlinear $\mu$-Function Crisis:\\
Can mochi\_class Handle DFD?\\[4pt]
{\normalsize Linear Horndeski Fails, Nonlinear AQUAL Perturbation Theory Required,}\\
{\normalsize Minimal CLASS Modifications, Observable Predictions Under Both Scenarios}}
\author{DFD Research Programme --- P10 Agent 9, Wave 4 Iteration 20 (Opus 4.6)}
\date{20 March 2026}

\begin{document}
\maketitle

%=============================================================================
\section*{W4i20 TOP 5 FINDINGS --- READ FIRST}
%=============================================================================

\begin{tcolorbox}[colback=blue!5!white, colframe=blue!60!black,
  title=\textbf{W4i20 TOP 5 FINDINGS --- RANKED BY IMPACT}]
\begin{enumerate}[nosep]

\item \textbf{[FINDING 1] EXISTING mochi\_class CANNOT HANDLE DFD.
  The Horndeski $\alpha$-parametrization is intrinsically LINEAR
  perturbation theory. DFD's cosmological perturbations live in the
  deep-MOND regime ($\mu(x) \approx x$, $x \ll 1$), where the field
  equation is NONLINEAR in the perturbation amplitude. The standard
  Horndeski QSA gives $\Geff = G_N$ (Agent 4's result), which is
  correct for the LINEAR theory but misses the MOND enhancement
  entirely. A parameter file for mochi\_class will reproduce
  Agent 4's catastrophic $P(k)$ deficit, not Agent 12's rescue.}\\
  \textbf{Source: Agent 4 (W4i19\_P4\_update, Eq.~(18)):
  $\Geff = G_N$ from strict QSA of linearized equation.
  Agent 12 (W4i19\_Cosmo\_update, Part II): deep-MOND requires
  nonlinear $\mu(x) \approx x$, giving $\Geff \propto 1/\sqrt{\delta}$.
  mochi\_class (Cataneo et al.\ 2024, arXiv:2407.11968): linear
  Horndeski perturbation theory only.}\\
  \textbf{Impact: EXISTENTIAL. The i19 ``zero new code'' piggyback
  strategy is DEAD. Custom nonlinear code is required.}

\item \textbf{[FINDING 2] THE FUNDAMENTAL DISAGREEMENT IS RESOLVED:
  BOTH AGENTS ARE CORRECT IN THEIR REGIMES.
  Agent 4's $\Geff = G_N$ is correct for the LINEARIZED perturbation
  equation (expand around $\nabla\bar\psi = 0$, keep first order in
  $\delta$). Agent 12's $\Geff \gg G_N$ is correct for the NONLINEAR
  quasi-static equation (AQUAL), where $\mu(|\nabla\psi|/\astar)$ with
  $\mu(0) = 0$ makes the linearization DEGENERATE. The physical question
  is: which regime applies to cosmological perturbations? Answer:
  the deep-MOND (nonlinear) regime, because cosmological accelerations
  $g_N \ll a_0$ on all observable scales (Agent 12, W4i19\_Cosmo\_update,
  \S Step 5: $y_N = g_N/a_0 \approx 0.05$ at recombination for
  $k = 0.1$~Mpc$^{-1}$). The linearized treatment is self-inconsistent:
  it assumes $\mu \approx 1$ (Newtonian background) but this requires
  $|\nabla\psi|/\astar \gg 1$, which is FALSE for cosmological
  perturbations.}\\
  \textbf{Source: Agent 4 (W4i19\_P4\_update, \S3.4, Eq.~(27)):
  $\Geff = G_N$ requires Newtonian background.
  Agent 12 (W4i19\_Cosmo\_update, \S Step 3--5): $y_N \ll 1$ on all
  observable scales, confirming deep-MOND.}\\
  \textbf{Impact: CRITICAL. Settles the disagreement in favor of
  Agent 12's deep-MOND scenario.}

\item \textbf{[FINDING 3] MINIMAL CLASS MODIFICATION FOR NONLINEAR
  DFD: $\sim$250 LINES OF C, BUT WITH A FUNDAMENTAL COMPLICATION.
  The deep-MOND $\Geff(k,\delta,z) = G_N\sqrt{a_0 k/(G\bar\rho_b
  \delta_b a)}$ depends on the perturbation AMPLITUDE $\delta_b$,
  making the growth equation nonlinear: $\ddot\delta + 2H\dot\delta
  \propto \sqrt{\delta}$. Standard Boltzmann codes (CLASS, CAMB)
  solve LINEAR ODEs. Implementing this requires either:
  (a) iterative solver (solve linear, update $\Geff$ with solution,
  repeat until convergence), or (b) direct nonlinear ODE integration
  (replace CLASS's linear growth ODE with RK4/RK45 for the nonlinear
  equation). Option (b) is more robust and adds $\sim$250 lines to
  CLASS. Neither mochi\_class nor any existing MG Boltzmann code
  supports amplitude-dependent $\Geff$.}\\
  \textbf{Source: Agent 12 (W4i19\_Cosmo\_update, Eq.~(F2.4)):
  $\Geff \propto 1/\sqrt{\delta_b}$. CLASS source code:
  perturbations.c uses linear ODE system.}\\
  \textbf{Impact: HIGH. Defines the actual coding task. 2--3
  developer-weeks, not 1.}

\item \textbf{[FINDING 4] OBSERVABLE PREDICTIONS UNDER BOTH SCENARIOS.\\
  LINEARIZED ($\Geff = G_N$, Agent 4):
  $\sigma_8 \approx 0.03$, $P(k)/P^{\rm CDM}(k) \sim 10^{-6}$
  at $k > 0.1$~Mpc$^{-1}$. CATASTROPHIC.
  $C_\ell^{TT}$: identical to baryon-only $\Lambda$CDM (no CDM).
  Ruled out by all LSS data.\\
  DEEP-MOND ($\Geff \gg G_N$, Agent 12):
  $\sigma_8 \approx 0.5$--$2$ (uncertain due to Silk damping and
  MOND$\to$Newton transition).
  $P(k)/P^{\rm CDM}(k) \sim 0.5$--$2$ at $k \sim 0.1$~Mpc$^{-1}$
  (VIABLE). Excess large-scale power at $k < 0.01$~Mpc$^{-1}$
  (potentially testable). $C_\ell^{TT}$: enhanced ISW at low $\ell$
  from scale-dependent growth, but magnitude uncertain without
  numerical integration.\\
  The code is needed to distinguish these scenarios numerically.}\\
  \textbf{Source: Agent 4 (W4i19\_P4\_update, \S3.4): linearized
  $\sigma_8 = 0.03$. Agent 12 (W4i19\_Cosmo\_update, Part III):
  deep-MOND $\delta_b(0) \sim 12$ from $a^2$ growth.}\\
  \textbf{Impact: EXISTENTIAL. The two scenarios predict different
  universes.}

\item \textbf{[FINDING 5] NO EXISTING BOLTZMANN CODE IMPLEMENTS
  MOND-LIKE NONLINEAR PERTURBATION THEORY. The Skordis-Zlosnik AeST
  code (used for their 2021 PRL CMB fit, arXiv:2007.00082) is NOT
  publicly available. CLASS-PT (Ivanov et al.\ 2020) handles
  nonlinear PT for CDM but not modified gravity with
  amplitude-dependent $\Geff$. MGCAMB and EFTCAMB use linear EFT.
  mochi\_class uses linear Horndeski. The DFD nonlinear perturbation
  code would be the FIRST public implementation of amplitude-dependent
  $\Geff$ in a Boltzmann solver.}\\
  \textbf{Source: Web search (2026-03-20). Skordis \& Zlosnik (2021,
  PRL 127, 161302): CMB fit shown but code not released. mochi\_class
  (Cataneo et al.\ 2024, arXiv:2407.11968): linear Horndeski only.
  CLASS-PT (Ivanov et al.\ 2020, arXiv:2004.10607): one-loop CDM
  only.}\\
  \textbf{Impact: HIGH. DFD requires a genuinely novel Boltzmann
  code extension. Also an opportunity: first-of-its-kind code.}

\end{enumerate}
\end{tcolorbox}

\begin{tcolorbox}[colback=red!5!white, colframe=red!60!black,
  title=\textbf{INCONSISTENCIES AND FLAGS}]
\begin{enumerate}[nosep]
\item \textbf{FLAG (CRITICAL)}: Finding 2 resolves the Agent 4 vs
  Agent 12 disagreement in favor of the deep-MOND scenario.
  However, this resolution ASSUMES that the cosmological background
  has $\nabla\bar\psi = 0$ (homogeneous). If there exists a nonzero
  cosmological gradient (e.g., from the Hubble flow contribution to
  the DFD field), the linearization point changes and the deep-MOND
  conclusion may be modified. The W4i19\_Cosmo\_update (Part I, \S
  ``Extending to the Nonlinear Field Equation'') identifies this
  issue: the linearization around $|\nabla\psi| = 0$ is degenerate
  for $\mu(0) = 0$. Resolution requires understanding the cosmological
  background psi-field, which depends on the DFD background equations.
  Since DFD does NOT modify $H(z)$, the background $\psi$ may indeed
  be trivial ($\nabla\bar\psi = 0$), but this needs verification from
  the action.

\item \textbf{FLAG}: The deep-MOND growth $\delta \propto a^2$
  gives $\delta_b(0) \sim 12$ starting from $\delta_b(z_{\rm dec})
  \sim 10^{-5}$. This estimate ignores: (a) Silk damping, which
  suppresses $\delta_b$ at $k > 0.1$~Mpc$^{-1}$ by factors of
  10--100; (b) the baryon-photon coupling before $z_{\rm dec}$, which
  prevents baryon growth; (c) the transition from deep-MOND to
  Newtonian as $\delta$ grows. Including Silk damping reduces the
  initial amplitude by $\sim e^{-k^2/k_{\rm Silk}^2}$, potentially
  re-introducing a deficit on small scales.

\item \textbf{FLAG}: The i19 P10 $\Geff$ formula (Eq.~(F1.1) of
  W4i19\_P10\_update) uses a different form from Agent 12's deep-MOND
  result. The P10 formula has $\Geff/G_N = 1/[\mu(x_0) + (k/k_H)^2
  x_0 \mu'(x_0)]$, which was derived from a linearized treatment
  with background gradient $x_0 = \astar a/k$. This formula gives
  $\Geff/G_N \approx 1/x_0 = k/(\astar a)$ for small $x_0$, which
  is DIFFERENT from Agent 12's $\Geff/G_N = \sqrt{a_0 k/(G\bar\rho_b
  \delta a)}$. The difference arises because the P10 formula linearizes
  around a fictitious ``background gradient'' while Agent 12 solves the
  full AQUAL algebraic relation. Agent 12's treatment is more
  physically motivated but requires verification.

\item \textbf{CORRECTION}: The i19 P10 update stated that mochi\_class
  MVP might be achievable with ``ZERO new code --- just a parameter
  file.'' This is INCORRECT. Even the linearized scenario requires
  understanding what $\alpha_K(a) \approx 0.014$ does in mochi\_class
  with all other $\alpha$'s zero --- this is a degenerate limit where
  the scalar field has no effect on perturbations in the QSA (because
  $\alpha_B = 0$ means no braiding, and kineticity alone does not
  modify the Poisson equation in the QSA). The parameter-file approach
  would give $\Geff = G_N$ identically, confirming Agent 4's result
  but not testing the deep-MOND scenario.
\end{enumerate}
\end{tcolorbox}

\newpage
\tableofcontents
\newpage

%=============================================================================
\section{Finding 1: Why Existing mochi\_class Cannot Handle DFD}
\label{sec:mochi-fails}
%=============================================================================

\subsection{The Linear Horndeski Framework}

The mochi\_class code (Cataneo et al.\ 2024, arXiv:2407.11968)
implements the Horndeski effective field theory of dark energy at
the level of \textbf{linear cosmological perturbation theory}. The
four $\alpha$-functions ($\alpha_K$, $\alpha_B$, $\alpha_M$,
$\alpha_T$) parameterize deviations from GR in the linearized
Einstein equations. The quasi-static approximation in mochi\_class
modifies the metric potentials $\Phi$ and $\Psi$ via:
\begin{equation}
-k^2\Phi = 4\pi G a^2 \mu(a,k)\,\bar\rho\,\delta, \qquad
-k^2(\Phi + \Psi) = 8\pi G a^2 \Sigma(a,k)\,\bar\rho\,\delta,
\label{eq:mochi-QSA}
\end{equation}
where $\mu(a,k)$ and $\Sigma(a,k)$ are time- and scale-dependent
functions determined by the $\alpha_i(a)$. These functions are
\textbf{independent of the perturbation amplitude $\delta$}.

\subsection{DFD's Mapping and Its Failure}

DFD maps onto $\alpha_T = \alpha_M = \alpha_B = 0$, $\alpha_K(a)
\neq 0$. In this subspace:
\begin{itemize}
\item $\alpha_T = 0$: $c_T = c$, no tensor speed anomaly.
  (Consistent with GW170817.)
\item $\alpha_M = 0$: No running Planck mass.
\item $\alpha_B = 0$: No kinetic braiding.
\item $\alpha_K(a) \approx 0.014$: Kineticity only.
\end{itemize}

In the Horndeski QSA with $\alpha_B = \alpha_M = 0$, the modified
Poisson equation parameters become (Bellini \& Sawicki 2014,
JCAP 1407, 050):
\begin{equation}
\mu_{\rm QSA}(a,k) = 1, \qquad \Sigma_{\rm QSA}(a,k) = 1.
\label{eq:horn-QSA-DFD}
\end{equation}

\textbf{Result}: In the Horndeski linear perturbation framework with
DFD's $\alpha$ values, $\Geff = G_N$ identically. The scalar field's
kineticity $\alpha_K$ affects the perturbation equations only through
the scalar sound speed and anisotropic stress, both of which vanish
in the QSA for $\alpha_B = 0$.

\textbf{This means}: Running mochi\_class with DFD's $\alpha$ values
reproduces Agent 4's catastrophic scenario ($\sigma_8 \approx 0.03$,
$P(k) \sim 10^{-6} P^{\rm CDM}$). The MOND enhancement is
\textbf{invisible} to the linear Horndeski framework.

\subsection{Why the MOND Enhancement is Nonlinear}

The DFD field equation in the quasi-static limit is the AQUAL equation
(W4i19\_Cosmo\_update, Eq.~(QS.6)):
\begin{equation}
\nabla\cdot\!\left[\mu\!\left(\frac{|\nabla\psi|}{\astar}\right)
\nabla\psi\right] = -\frac{8\pi G}{c^2}\rho_b.
\label{eq:AQUAL}
\end{equation}

With $\mu(x) = x/(1+x)$:
\begin{itemize}
\item $\mu(0) = 0$: The equation is \textbf{degenerate} at zero
  gradient. Linearization around $|\nabla\bar\psi| = 0$ gives
  $0 \cdot \nabla^2\delta\psi = \text{source}$, which is singular.
\item The physics requires solving the full nonlinear equation.
  For small perturbations where $|\nabla\delta\psi|/\astar \ll 1$
  (deep-MOND regime), $\mu \approx x$ and the equation becomes:
  \begin{equation}
  \nabla\cdot\!\left[\frac{|\nabla\psi|}{\astar}\,\nabla\psi\right]
  = -\frac{8\pi G}{c^2}\rho_b,
  \end{equation}
  which is quadratic in $\nabla\psi$ --- inherently nonlinear.
\end{itemize}

This is why no linear Boltzmann code can capture DFD's MOND sector:
the very thing that makes DFD interesting (the $\mu$-function
crossover with $\mu(0) = 0$) makes the perturbation theory nonlinear
at the lowest order.


%=============================================================================
\section{Finding 2: Resolution of the Agent 4 vs Agent 12 Disagreement}
\label{sec:resolution}
%=============================================================================

\subsection{Agent 4's Derivation (Linearized)}

Agent 4 (W4i19\_P4\_update) starts from the linearized wave equation:
\begin{equation}
\ddot\delta\psi_k + 3H\dot\delta\psi_k + \frac{c^2 k^2}{a^2}\delta\psi_k
= 8\pi G\,\bar\rho_b\,\delta_b.
\end{equation}

In the quasi-static limit ($k \gg aH/c$), this gives:
\begin{equation}
\frac{c^2 k^2}{a^2}\delta\psi_k = 8\pi G\,\bar\rho_b\,\delta_b,
\end{equation}
and substituting into the baryon growth equation:
\begin{equation}
\ddot\delta_b + 2H\dot\delta_b = 4\pi G\,\bar\rho_b\,\delta_b.
\end{equation}

This gives $\Geff = G_N$ exactly. The derivation is mathematically
correct \textbf{within the linearized theory}.

\subsection{Agent 12's Derivation (Nonlinear AQUAL)}

Agent 12 (W4i19\_Cosmo\_update) starts from the full quasi-static
AQUAL equation~(\ref{eq:AQUAL}) and derives the MOND algebraic
relation:
\begin{equation}
\mu(g/a_0)\cdot g = g_N, \qquad
g = g_N / \mu(g/a_0).
\end{equation}

For $g_N \ll a_0$ (deep-MOND): $\mu \approx g/a_0$, giving
$g = \sqrt{g_N \cdot a_0}$ and:
\begin{equation}
\Geff = \frac{g}{g_N}\,G_N = \sqrt{\frac{a_0}{g_N}}\,G_N \gg G_N.
\end{equation}

This is physically correct but operates in a regime where the
linearized treatment breaks down.

\subsection{Which Regime Applies?}

The key diagnostic is the MOND parameter $y_N = g_N/a_0$, where
$g_N$ is the Newtonian gravitational acceleration produced by the
perturbation. Agent 12 computed (W4i19\_Cosmo\_update, \S Step 5):

\begin{center}
\begin{tabular}{lccc}
\toprule
Epoch & $k$ (Mpc$^{-1}$) & $y_N = g_N/a_0$ & Regime \\
\midrule
$z = 1100$, $\delta = 10^{-5}$ & 0.01 & 0.5 & Crossover \\
$z = 1100$, $\delta = 10^{-5}$ & 0.1 & 0.05 & Deep MOND \\
$z = 1100$, $\delta = 10^{-5}$ & 1.0 & 0.005 & Deep MOND \\
$z = 0$, $\delta = 1$ & 0.1 & $4\ee{-4}$ & Deep MOND \\
\bottomrule
\end{tabular}
\end{center}

On ALL scales relevant for the matter power spectrum, cosmological
perturbations are in the deep-MOND regime. Agent 4's linearized
treatment assumes $\mu \approx 1$ (Newtonian), which requires
$y_N \gg 1$ --- this is satisfied only on sub-galactic scales with
large overdensities.

\textbf{Conclusion}: Agent 12's deep-MOND scenario is the physically
correct one for cosmological perturbation theory. Agent 4's result
is a valid mathematical limit but does not describe our universe.

\subsection{The Self-Inconsistency of the Linear Treatment}

The linearized DFD perturbation equation was derived by expanding
the action to second order around the background $\bar\psi = \text{const}$,
yielding $K''(\bar\Delta) = 1$ (from $K''(0) = 1$). This expansion
\textbf{implicitly assumes} that the perturbation $\delta\psi$ is small
compared to the nonlinear scale set by $\mu$. But $\mu(x) = x/(1+x)$
with $\mu(0) = 0$ means the function $K(\Delta)$ has $K'(0) = 0$
(the MOND condition). The expansion around $\Delta = 0$ gives a kinetic
term proportional to $K''(0)(\delta\Delta)^2$, but the actual field
equation involves $\mu(|\nabla\psi|/\astar) = K'(\Delta)/\astar$ which
vanishes at the background.

The linearized treatment effectively replaces $\mu(x) \to 1$ (its
asymptotic value), discarding the very feature ($\mu(0) = 0$) that
makes DFD a modified gravity theory. This is analogous to linearizing
the square-root in deep-MOND kinematics: one loses the $1/r$ force law
that is the whole point.


%=============================================================================
\section{Finding 3: Minimal CLASS Modification for Nonlinear DFD}
\label{sec:code}
%=============================================================================

\subsection{The Problem with Standard Boltzmann Codes}

All existing Einstein-Boltzmann solvers (CLASS, CAMB, and their
extensions) solve a SYSTEM OF LINEAR ODEs for the perturbation
variables $(\delta_\gamma, v_\gamma, \delta_b, v_b, \delta_{\rm CDM},
v_{\rm CDM}, \Phi, \Psi, \ldots)$. The linearity is essential for:
\begin{enumerate}
\item Mode-by-mode evolution (each $k$ evolves independently).
\item Transfer function formalism ($T(k) = \delta(k,z)/\delta_{\rm prim}(k)$).
\item Superposition (sum initial conditions for $C_\ell$ integration).
\end{enumerate}

DFD's deep-MOND perturbation theory breaks ALL three:
\begin{enumerate}
\item $\Geff$ depends on $\delta_b$, so different modes couple.
\item $T(k)$ is not a linear transfer function; it depends on
  the primordial amplitude.
\item Superposition fails: $\delta_1 + \delta_2$ does not give
  the sum of individual evolutions.
\end{enumerate}

\subsection{Strategy: Iterative Linearization}

The practical approach is \textbf{iterative linearization}:
\begin{enumerate}
\item Run CLASS with $\Geff = G_N$ (standard, iteration 0).
\item Extract $\delta_b(k,z)$ from iteration 0.
\item Compute $\Geff(k,z) = G_N \sqrt{a_0 k / (G\bar\rho_b
  |\delta_b(k,z)| a)}$ using the iteration-0 solution.
\item Re-run CLASS with the scale- and time-dependent $\Geff(k,z)$.
\item Repeat until $P(k)$ converges (typically 3--5 iterations).
\end{enumerate}

This preserves CLASS's linear ODE infrastructure while capturing the
nonlinear $\Geff$. Each iteration is a standard CLASS run with a
modified $\Geff(k,a)$ lookup table.

\subsection{Code Specification}

\begin{tcolorbox}[colback=green!5!white, colframe=green!60!black,
  title=\textbf{NONLINEAR DFD CLASS: Complete Coding Specification}]

\textbf{Architecture}: Wrapper script (Python) + modified CLASS.

\textbf{Step 1: Header file \texttt{include/dfd\_nonlinear.h}
($\sim$50 lines)}

\begin{lstlisting}[caption={dfd\_nonlinear.h --- header for nonlinear DFD},numbers=none]
#ifndef DFD_NONLINEAR_H
#define DFD_NONLINEAR_H

struct dfd_params {
  short  dfd_on;
  double a0;         /* MOND accel [m/s^2]           */
  double a_star;     /* 2*a0/c^2 in CLASS units      */
  int    iteration;  /* current iteration number      */
  /* Lookup table from previous iteration */
  int    n_k;
  int    n_z;
  double * k_table;  /* k values [1/Mpc]             */
  double * z_table;  /* z values                      */
  double * delta_b;  /* delta_b(k,z) from prev iter   */
  double * Geff_table; /* G_eff/G_N(k,z) precomputed */
};

/* Compute G_eff/G_N using lookup table */
double dfd_Geff_ratio(double k, double a,
                      struct dfd_params * pdfd);

/* Initialize lookup table from file */
int dfd_read_delta_table(char * filename,
                         struct dfd_params * pdfd);

/* Compute G_eff table from delta table */
int dfd_compute_Geff_table(struct dfd_params * pdfd);

/* Interpolate G_eff at arbitrary (k, z) */
double dfd_interpolate_Geff(double k, double z,
                            struct dfd_params * pdfd);

#endif
\end{lstlisting}

\textbf{Step 2: Implementation \texttt{source/dfd\_nonlinear.c}
($\sim$150 lines)}

\begin{lstlisting}[caption={dfd\_nonlinear.c --- core implementation},numbers=none]
#include "dfd_nonlinear.h"
#include <math.h>
#include <stdio.h>
#include <stdlib.h>

double dfd_Geff_ratio(double k, double a,
                      struct dfd_params * pdfd) {
  if (!pdfd->dfd_on) return 1.0;
  if (pdfd->iteration == 0) return 1.0; /* first iter: GR */

  /* Interpolate from precomputed table */
  double z = 1.0/a - 1.0;
  return dfd_interpolate_Geff(k, z, pdfd);
}

int dfd_compute_Geff_table(struct dfd_params * pdfd) {
  /* For each (k, z) point in the table:
   *   g_N = G * rho_bar(z) * |delta_b(k,z)| * a / k
   *   y_N = g_N / a0
   *   If y_N < 1 (deep-MOND):
   *     G_eff/G_N = sqrt(a0/g_N) = 1/sqrt(y_N)
   *   If y_N > 1 (Newtonian):
   *     G_eff/G_N = 1/mu(y_N) where mu(x)=x/(1+x)
   *   Full interpolation: solve mu(x)*x = y_N for x,
   *     then G_eff/G_N = x/y_N
   */
  int ik, iz;
  double c2 = 8.988e16; /* c^2 in m^2/s^2 */
  double rho_crit_0 = 8.62e-27; /* kg/m^3 */
  double Omega_b = 0.049;

  for (iz = 0; iz < pdfd->n_z; iz++) {
    double z = pdfd->z_table[iz];
    double a = 1.0 / (1.0 + z);
    double rho_b = Omega_b * rho_crit_0 * pow(1+z, 3);

    for (ik = 0; ik < pdfd->n_k; ik++) {
      double k = pdfd->k_table[ik];
      double delta = fabs(pdfd->delta_b[iz*pdfd->n_k+ik]);
      if (delta < 1e-15) delta = 1e-15; /* floor */

      /* Newtonian acceleration of perturbation */
      /* g_N ~ G * rho_b * delta * (2*pi*a / k) */
      /* in physical coords */
      double k_phys = k / a; /* physical k */
      double lambda = 2*M_PI / k_phys; /* physical lambda */
      double g_N = 6.674e-11 * rho_b * delta * lambda;
      double y_N = g_N / pdfd->a0;

      /* Solve mu(x)*x = y_N for the simple mu */
      /* mu(x) = x/(1+x), so x^2/(1+x) = y_N */
      /* x^2 - y_N*x - y_N = 0 */
      /* x = (y_N + sqrt(y_N^2 + 4*y_N)) / 2 */
      double disc = y_N*y_N + 4.0*y_N;
      double x = (y_N + sqrt(disc)) / 2.0;

      /* G_eff/G_N = x / y_N */
      double ratio;
      if (y_N > 1e-20) {
        ratio = x / y_N;
      } else {
        ratio = 100.0; /* cap */
      }

      /* Safety clamp */
      if (ratio < 1.0) ratio = 1.0;
      if (ratio > 200.0) ratio = 200.0;

      pdfd->Geff_table[iz*pdfd->n_k+ik] = ratio;
    }
  }
  return 0; /* success */
}
\end{lstlisting}

\textbf{Step 3: Hook into \texttt{source/perturbations.c}
($\sim$25 lines)}

In the function \texttt{perturbations\_einstein()}, after the Poisson
equation evaluation (the line computing
\texttt{pvecmetric[ppw->index\_mt\_phi\_prime]}):

\begin{lstlisting}[numbers=none]
#ifdef DFD_MODULE
if (ppr->dfd.dfd_on && ppr->dfd.iteration > 0) {
  double g_ratio = dfd_Geff_ratio(
    k,
    pvecback[pba->index_bg_a],
    &ppr->dfd);
  /* Scale the gravitational source terms */
  /* This modifies the Poisson equation:
   *   k^2 Phi = -4*pi*G_eff*a^2*rho*delta
   * by multiplying the RHS by g_ratio */

  /* Modify the metric perturbation source */
  pvecmetric[ppw->index_mt_phi_prime] *=
    g_ratio;
  /* Also modify psi for the lensing potential */
  pvecmetric[ppw->index_mt_psi] *= g_ratio;
}
#endif
\end{lstlisting}

\textbf{Step 4: Python iteration wrapper ($\sim$80 lines)}

\begin{lstlisting}[language=Python,numbers=none,
  caption={iterate\_dfd.py --- convergence loop}]
import subprocess, numpy as np

def run_class_iteration(iteration, delta_file=None):
    """Run CLASS with DFD module at given iteration."""
    params = {
        'dfd_on': 'yes',
        'dfd_a0': '1.2e-10',
        'dfd_iteration': str(iteration),
    }
    if delta_file:
        params['dfd_delta_file'] = delta_file
    # Write .ini file and run CLASS
    # ... (standard CLASS parameter file generation)
    subprocess.run(['./class', 'dfd_iter.ini'])

def extract_delta_b(output_dir):
    """Read delta_b(k,z) from CLASS output."""
    # Read transfer functions at multiple z
    # Return (k_array, z_array, delta_b_array)
    pass

def check_convergence(Pk_old, Pk_new, tol=0.01):
    """Check if P(k) has converged."""
    ratio = Pk_new / Pk_old
    return np.all(np.abs(ratio - 1) < tol)

# Main iteration loop
max_iter = 10
for i in range(max_iter):
    if i == 0:
        run_class_iteration(0)
    else:
        run_class_iteration(i, f'delta_iter{i-1}.dat')
    delta_b = extract_delta_b('output/')
    np.savetxt(f'delta_iter{i}.dat', delta_b)
    if i > 0 and check_convergence(Pk_old, Pk_new):
        print(f"Converged at iteration {i}")
        break
    Pk_old = Pk_new
\end{lstlisting}

\textbf{Total new code}: $\sim$250 lines C + $\sim$80 lines Python.

\textbf{Timeline}: 2--3 developer-weeks for a CLASS expert (additional
week vs.\ i19 MVP due to iteration infrastructure and convergence
testing).
\end{tcolorbox}


%=============================================================================
\section{Finding 4: Observable Predictions Under Both Scenarios}
\label{sec:predictions}
%=============================================================================

\subsection{Scenario A: Linearized ($\Geff = G_N$)}

This is what mochi\_class or the i19 linear MVP would produce.
DFD reduces to a baryon-only universe with standard Newtonian gravity:

\begin{center}
\begin{tabular}{lcc}
\toprule
Observable & DFD (linearized) & \LCDM \\
\midrule
$\sigma_8$ & $0.03 \pm 0.01$ & $0.811 \pm 0.006$ \\
$P(k=0.1)/P^{\rm CDM}(k=0.1)$ & $\sim 10^{-6}$ & 1 \\
$C_\ell^{TT}$ (low $\ell$) & Baryon-only ISW & Planck fit \\
$f\sigma_8(z=0.5)$ & $\sim 0.02$ & $0.47 \pm 0.04$ \\
\bottomrule
\end{tabular}
\end{center}

\textbf{Derivation of $\sigma_8 \approx 0.03$} (from W4i19\_Cosmo\_update,
Finding 3):

Baryons start growing after decoupling at $z = 1100$ with
$\delta_b \sim 10^{-5}$. With $\Geff = G_N$ and no CDM potential
wells, growth follows $\delta_b \propto a$ (matter era) then
stalls ($\Lambda$ era). Total growth factor $\sim 1100 \times 0.8
\approx 900$. Final $\delta_b \sim 10^{-5} \times 900 \sim 10^{-2}$.

$\sigma_8 \propto \delta_b(0) \Rightarrow \sigma_8^{\rm DFD}
\approx 0.03$. This is 27x below the observed value.

\textbf{Verdict}: CATASTROPHICALLY excluded.

\subsection{Scenario B: Deep-MOND ($\Geff \gg G_N$)}

This is what the nonlinear iterative code would produce:

\begin{center}
\begin{tabular}{lcc}
\toprule
Observable & DFD (deep-MOND) & \LCDM \\
\midrule
$\sigma_8$ & $0.5$--$2$ (uncertain) & $0.811 \pm 0.006$ \\
$P(k=0.1)/P^{\rm CDM}(k=0.1)$ & $\sim 0.5$--$2$ & 1 \\
$C_\ell^{TT}$ (low $\ell$) & Enhanced ISW & Planck fit \\
$f\sigma_8(z=0.5)$ & $\sim 0.3$--$0.8$ & $0.47 \pm 0.04$ \\
\bottomrule
\end{tabular}
\end{center}

\textbf{Derivation of $\sigma_8 \approx 0.5$--$2$} (from
W4i19\_Cosmo\_update, Finding 3):

Deep-MOND growth: $\delta_b \propto a^2$. From $z = 1100$:
$\delta_b(0) \sim 10^{-5} \times (1101)^2 \sim 12$. But Silk
damping reduces the initial amplitude at $k \sim 0.1$~Mpc$^{-1}$
by a factor $\sim e^{-(k/k_{\rm Silk})^2}$ with $k_{\rm Silk}
\approx 0.07$~Mpc$^{-1}$. For $k = 0.1$: damping factor
$\sim e^{-2} \approx 0.14$. Effective $\delta_b(0) \sim 1.7$.

Including the Newtonian transition (which slows growth when
$\delta > 1$): $\sigma_8 \sim 0.5$--$1$.

If Silk damping is stronger (as in a baryon-only universe where
photon-baryon coupling is tighter), the reduction could push
$\sigma_8$ lower, toward $0.3$--$0.5$.

\textbf{Verdict}: VIABLE but requires numerical confirmation.
The $S_8$ tension ($S_8^{\rm Planck} = 0.83$, $S_8^{\rm KiDS}
= 0.76$) might actually prefer a slightly lower $\sigma_8$.

\subsection{$C_\ell^{TT}$ Predictions}

\textbf{Scenario A}: Standard baryon-only CMB. The acoustic peaks
are at the correct positions (DFD does not modify the sound horizon).
However, the peak height ratios change: the odd-even asymmetry
(driven by baryon loading) is maximal (no CDM). The third peak is
suppressed relative to the second. This disagrees with Planck at
$> 10\sigma$ because Planck's third peak measurement constrains
$\Omega_{\rm CDM}$.

\textbf{Scenario B}: The pre-decoupling physics may be modified
if $\Geff \neq G_N$ during the radiation era. However, in the
deep-MOND regime, the radiation-era perturbations (which are large:
$\delta\rho_\gamma/\bar\rho_\gamma \sim 10^{-5}$ but $\bar\rho$
is huge) may have $y_N \gg 1$, placing them in the Newtonian regime.
If so, the CMB acoustic oscillations are unchanged, and only the
late-time ISW effect is modified by the enhanced growth.

\textbf{This is a critical test}: the CMB peak structure constrains
the pre-decoupling gravity. If deep-MOND applies during the acoustic
oscillations, the CMB power spectrum would be drastically altered.
The Skordis-Zlosnik AeST theory avoids this by having a separate
cosmological mechanism (the aether field acts as dark matter for
CMB purposes). DFD would need a similar mechanism --- or the
radiation-era perturbations must be in the Newtonian regime.

\subsection{$P(k)$ Shape}

The deep-MOND growth $\delta \propto a^2$ is \textbf{scale-independent}
in the power-law regime. This means:
\begin{equation}
P_{\rm DFD}(k) = P_{\rm initial}(k) \times T_{\rm Silk}^2(k)
\times D_{\rm MOND}^2(z),
\end{equation}
where $D_{\rm MOND}(z) \propto a^2(z)$ and $T_{\rm Silk}(k)$ is the
Silk damping transfer function. The SHAPE of $P(k)$ is determined
by the primordial spectrum and Silk damping only --- there is no CDM
``turnover'' at $k_{\rm eq}$.

In \LCDM, the $P(k)$ turnover at $k_{\rm eq} \approx 0.01$~Mpc$^{-1}$
is a signature of the matter-radiation equality epoch. In DFD (Scenario B),
this turnover is ABSENT (or modified). The observed $P(k)$ turnover at
$k \sim 0.01$--$0.02$~Mpc$^{-1}$ is one of the strongest constraints
on dark matter. DFD's Scenario B must reproduce this feature through
a different mechanism (e.g., the MOND-to-Newton transition providing
a scale), or it faces a serious tension.


%=============================================================================
\section{Finding 5: Literature on MOND Boltzmann Codes}
\label{sec:literature}
%=============================================================================

\subsection{Skordis-Zlosnik AeST}

The Aether Scalar Tensor (AeST) theory of Skordis \& Zlosnik (2021,
PRL 127, 161302, arXiv:2007.00082) is the most successful relativistic
MOND theory to date. Key features:
\begin{itemize}
\item Reproduces CMB $C_\ell^{TT}$ and matter $P(k)$ consistent with
  Planck + BOSS.
\item Has a unit-timelike vector field (aether) that provides an
  effective dark matter component at early times.
\item Reduces to MOND in the quasi-static weak-field limit.
\item The Boltzmann code used for their CMB fit is \textbf{NOT publicly
  available}. Web search (2026-03-20) found no GitHub repository or
  public release.
\item The Hamiltonian formalism was developed in Skordis et al.\ (2024,
  PRD 110, 044015, arXiv:2307.15126), showing 6 physical degrees of
  freedom.
\end{itemize}

DFD differs from AeST in that DFD has no aether field --- it uses a
single scalar $\psi$ on a fixed flat background. This means DFD lacks
AeST's early-time dark matter analog, which is what allows AeST to
fit the CMB peaks. DFD's CMB fit strategy must be different (either
the $\mu$-function modifies the acoustic oscillations, or the theory
accepts the CDM-free CMB as a prediction to be tested).

\subsection{Existing Modified Gravity Boltzmann Codes}

\begin{center}
\begin{tabular}{llll}
\toprule
Code & Framework & Perturbation theory & DFD compatible? \\
\midrule
mochi\_class & Horndeski $\alpha_i$ & Linear & No \\
hi\_class & Horndeski $\alpha_i$ & Linear & No \\
MGCAMB & $\mu$-$\Sigma$ & Linear & No \\
EFTCAMB & EFT of DE & Linear & No \\
CLASS-PT & CDM 1-loop PT & Nonlinear (CDM) & No \\
AeST code & AeST theory & Linear (custom) & Not public \\
\textbf{DFD-CLASS} & \textbf{AQUAL + iter.} & \textbf{Nonlinear} & \textbf{To build} \\
\bottomrule
\end{tabular}
\end{center}

\subsection{Closest Existing Work}

The closest existing implementation to what DFD needs is:
\begin{enumerate}
\item \textbf{MOND N-body simulations} (e.g., Phantom of RAMSES,
  Lughausen et al.\ 2015): These solve the full nonlinear AQUAL
  equation in an N-body context but do not produce CMB observables.
\item \textbf{The Skordis-Zlosnik private code}: This handles the
  full AeST perturbation theory including CMB, but is not public
  and uses a different theory.
\item \textbf{COPPER} (Noller 2022, github.com/noller/copper):
  A general parametrized perturbation code, but limited to linear theory.
\end{enumerate}


%=============================================================================
\section{Synthesis and Recommended Path Forward}
\label{sec:synthesis}
%=============================================================================

\subsection{The Existential Fork}

DFD's cosmological viability now depends on a single question:

\begin{tcolorbox}[colback=yellow!10!white, colframe=yellow!60!black,
  title=\textbf{THE QUESTION}]
\textbf{Does the deep-MOND enhanced growth ($\delta \propto a^2$)
produce a matter power spectrum $P(k)$ that matches observations,
after properly accounting for Silk damping, the MOND-to-Newton
transition, and the CMB peak structure?}

This question can ONLY be answered by the nonlinear iterative
DFD-CLASS code described in Finding 3.
\end{tcolorbox}

\subsection{Why mochi\_class Cannot Answer This}

To summarize the argument:
\begin{enumerate}
\item mochi\_class implements linear Horndeski perturbation theory.
\item DFD maps onto $\alpha_B = \alpha_M = \alpha_T = 0$, $\alpha_K
  \neq 0$.
\item In this subspace, the linear QSA gives $\Geff = G_N$ identically.
\item The MOND enhancement is a NONLINEAR effect from $\mu(0) = 0$.
\item No linear MG Boltzmann code captures this effect.
\item Running mochi\_class with DFD parameters would confirm the
  catastrophic scenario (Scenario A), which is the WRONG physical
  scenario.
\end{enumerate}

\subsection{Recommended Coding Priority}

\begin{enumerate}
\item \textbf{Phase 1 (1 week)}: Build the iterative DFD-CLASS with
  the deep-MOND $\Geff$ formula. Run 5 iterations. Check convergence.
  Produce $P(k)$ at $z = 0, 0.5, 1.0$.

\item \textbf{Phase 2 (1 week)}: Add CMB integration ($C_\ell^{TT}$).
  Check whether the enhanced ISW at low $\ell$ is within Planck
  uncertainties. Check acoustic peak structure.

\item \textbf{Phase 3 (1 week)}: Compare to BOSS/DESI $P(k)$ and
  $f\sigma_8$ measurements. Determine if DFD is viable, marginal,
  or excluded.

\item \textbf{Phase 4 (if viable)}: Full MCMC parameter estimation
  to constrain $a_0$ from cosmological data.
\end{enumerate}

\textbf{Estimated total cost}: 3--4 developer-weeks, \$15--20K.
This is 3x the i19 MVP estimate but is the MINIMUM for a
scientifically meaningful answer.


%=============================================================================
\section*{Cross-References and Provenance}
%=============================================================================

All claims in this document trace to:
\begin{itemize}
\item \textbf{Agent 4 (W4i19\_P4\_update)}: Linearized $\Geff = G_N$,
  QSA derivation, Eqs.~(18), (27), \S3.4.
\item \textbf{Agent 12 (W4i19\_Cosmo\_update)}: Deep-MOND regime,
  $\Geff \gg G_N$, $y_N$ calculation, $\delta \propto a^2$ growth,
  Parts I--III.
\item \textbf{Agent 9 (W4i19\_P10\_update)}: mochi\_class architecture,
  Horndeski mapping, $\alpha_K$ derivation, \S\S1--2.
\item \textbf{mochi\_class paper}: Cataneo et al.\ (2024),
  arXiv:2407.11968, linear Horndeski QSA framework.
\item \textbf{Skordis-Zlosnik}: PRL 127, 161302 (2021),
  arXiv:2007.00082; Hamiltonian formalism: PRD 110, 044015 (2024),
  arXiv:2307.15126.
\item \textbf{CLASS-PT}: Ivanov et al.\ (2020), arXiv:2004.10607.
\item \textbf{Web search}: Conducted 2026-03-20 for MOND Boltzmann
  codes, AeST code availability, mochi\_class limitations.
\end{itemize}

\end{document}
